% !TEX program = lualatex
% =====================================================================
% despacho-pesquisa.tex -- documento avulso, em pe:
%   levantamento de literatura e de fonte setorial sobre despacho de energia
%
% PROCEDENCIA. Levantamento de 06/09/2026, feito por busca na web e por
% consulta a API, e NAO por leitura do estoque de dado deste repositorio.
% Nenhum numero daqui foi medido em arquivo local: todos sao citados, e a
% distincao entre o que veio de fonte primaria e o que veio de resumo de
% segunda mao esta declarada na secao de procedencia, ao fim.
%
% REFERENCIAS conferidas campo a campo contra o registro do Crossref e da
% API do arXiv na mesma data -- autor, container, volume, pagina, ano e
% DOI. Onde o Crossref devolveu a versao de conferencia E a de periodico,
% a de periodico prevaleceu (`maceira2015` e `soares2017`; a primeira tem
% versao PSCC 2014, a segunda IEEE PES GM 2014). `chen2024` entra pela
% versao do IEEE Transactions on Power Systems, e nao pelo pre-print.
%
% O QUE NAO FOI CONFERIDO, e a objecao final diz:
%   - o texto da Lei 15.269/2025. Duas fontes secundarias a citam para
%     coisas distintas -- marco de modernizacao e acerto de curtailment --
%     e nao se abriu o texto para separar os dispositivos.
%   - o artigo de `gorka2024` por inteiro: a ACM devolveu HTTP 403. Metodo
%     e alvo saem do resumo e do registro, nao do corpo do artigo.
%   - a nota tecnica do ONS de 23/09/2025 sobre excedentes, so citada.
%
% Nao entra no `reporte.tex`. Compila sozinho:
%   latexmk -pdf despacho-pesquisa.tex
%
% Em pe, e nao deitado como `tarefas-dados.tex`: aqui o conteudo e prosa
% com duas tabelas estreitas.
%
% MARGEM de 1 cm, pedida em 06/09/2026. Isso poe a mancha em 19 cm, contra
% os 12,5 cm que a folha de estilo declara para documento de estudo. A
% medida de linha vai a ~150 caracteres, acima da faixa de 45 a 60 que o
% cabecalho da folha de estilo justifica. E escolha declarada, nao
% descuido: quem quiser a medida da folha volta ao geometry anterior,
% `top=2.4cm,bottom=2.4cm,textwidth=12.5cm,centering`.
% =====================================================================

% A folha de estilo pressupoe `book`: usa `\chaptertitlename` e os
% comandos de capitulo do `tocloft`. Em `article` ela nao carrega.
\documentclass[10pt,a4paper,oneside]{book}

\usepackage{iftex}
\ifPDFTeX\usepackage[utf8]{inputenc}\fi
\usepackage[brazilian]{babel}
\usepackage[a4paper,margin=1cm]{geometry}
\usepackage{estilo-reporte}

\pagestyle{plain}

\makeatletter
\renewenvironment{thebibliography}[1]
  {\list{\@biblabel{\@arabic\c@enumiv}}%
     {\settowidth\labelwidth{\@biblabel{#1}}%
      \leftmargin\labelwidth \advance\leftmargin\labelsep
      \usecounter{enumiv}\let\p@enumiv\@empty
      \renewcommand\theenumiv{\@arabic\c@enumiv}}%
   \sloppy\clubpenalty4000 \@clubpenalty\clubpenalty
   \widowpenalty4000 \sfcode`\.\@m}
  {\endlist}
\makeatother

\begin{document}

{\sffamily\bfseries\LARGE Despacho de energia: o problema, o caso brasileiro e o que o aprendizado de máquina faz com ele\par}
\vspace{0.4em}
{\sffamily\fontsize{8}{10}\selectfont\color{tinta2}%
Levantamento de literatura e de fonte setorial, 6 de setembro de 2026.
Nenhum número deste documento foi medido em arquivo local: todos são citados.\par}
\vspace{1.4em}

Despacho é um só problema de otimização resolvido em várias escalas de tempo;
o que varia entre países é \termo{quem decide e com que informação}. O Brasil
está no meio da transição que muda exatamente isso --- de despacho centralizado
por custo declarado para despacho por oferta --- e é essa transição, e não a
técnica de otimização, que redefine o que ``despacho'' significa como objeto de
estudo nos próximos anos. Quem ajusta modelo sobre a série atual está ajustando
sobre um regime com data de validade conhecida.

\section*{A formulação, e por que ela se parte em camadas}

O problema é um programa inteiro misto: comprometimento de unidades --- variável
binária, ligar ou não --- despacho econômico --- variável contínua, quanto ---
e restrições de rede sobrepostas às duas. Resolver tudo junto no tamanho real é
caro, e a prática mundial fatia por horizonte.

\vspace{0.4em}
\tabfont
\begin{tabular}{@{}
    >{\rag\arraybackslash}p{3.6cm}
    >{\rag\arraybackslash}p{5.6cm}
    >{\rag\arraybackslash}p{2.6cm}@{}}
\toprule
\cab{Camada} & \cab{O que decide} & \cab{Cadência} \\
\midrule
SCUC --- comprometimento com segurança &
quais unidades ligam, com custo de partida, tempo mínimo ligado e desligado &
dia anterior \\
\addlinespace[.4ex]
SCED --- despacho econômico com segurança &
quanto cada unidade gera, dado o comprometimento já fixado &
5 minutos em MISO e PJM \\
\bottomrule
\end{tabular}

\notas{O corte entre as duas camadas é convenção operacional, não propriedade
do problema. É ele que cria a janela em que a decisão de despacho já está
tomada antes de o recurso variável se manifestar --- e é dessa janela que sai
todo o corte de geração programado.}

\vspace{0.6em}
A dificuldade computacional se concentra no binário do SCUC, que não relaxa
bem, enquanto o SCED tem prazo de segundos e roda a cada cinco minutos nos
operadores americanos~\cite{chen2022}. As duas restrições --- combinatória de
um lado, prazo do outro --- explicam por que a literatura de aceleração se
divide em duas linhas que quase não conversam.

\section*{O caso brasileiro: preço é sombra de modelo, não oferta}

O SIN não despacha por oferta. Uma cadeia de modelos do CEPEL calcula o
despacho, e do dual dela sai o preço: NEWAVE no médio prazo, que produz a
função de custo futuro --- o valor da água; DECOMP no curto prazo, semanal; e
DESSEM no curtíssimo prazo, oficial para a programação do despacho desde 2020 e
para o PLD horário desde 2021~\cite{cepel2026,ccee2026}.

Três consequências distinguem o Brasil de qualquer mercado \emph{bid-based}, e
as três importam para quem modela.

\campo{O preço é um multiplicador de Lagrange, não uma oferta}
O CMO e o PLD são o custo marginal do modelo, não o que alguém pediu para
gerar. Isso torna o preço reproduzível em princípio e opaco na prática: ele
depende de uma função de custo futuro reestimada periodicamente, cujas revisões
a CCEE publica como evento próprio, com data de vigência.

\campo{A aversão ao risco é parâmetro, não resultado}
O CVaR entrou oficialmente em setembro de 2013 na operação, na formação do
preço e no planejamento da expansão. A abordagem por combinação convexa de
esperança e CVaR produz soluções mais conservadoras sem sacrificar custo
total~\cite{maceira2015}. O grau de conservadorismo, porém, é escolha
regulatória embutida no modelo que forma o preço --- não uma propriedade do
sistema físico.

\campo{A variabilidade do dual é defeito conhecido e não resolvido}
Modelos autorregressivos de vazão induzem variabilidade indesejada da solução
primal --- geração térmica --- e da dual --- custo marginal e preço
\emph{spot} --- altamente sensível às condições iniciais de
afluência~\cite{soares2017}. \termo{Quem prevê CMO está prevendo, em parte, a
sensibilidade numérica de um SDDP}, e não só a hidrologia. Um modelo estatístico
que acerte essa componente não descobriu física nenhuma.

\section*{A reforma que muda o estimando}

Em março de 2026 o MME submeteu a consulta pública uma proposta de
aperfeiçoamento das diretrizes para adoção da contabilização dual no Mercado de
Curto Prazo e de transição para um sistema baseado em ofertas de quantidade de
energia, a serem consideradas na otimização e na formação de preço de curto
prazo~\cite{cp218}. A agenda de 2026 traz também a modernização do marco
regulatório atribuída à Lei 15.269/2025~\cite{cp218,tagd2026}.

A consequência para modelagem é direta e pouco comentada. \termo{Num sistema por
oferta, prever despacho vira prever comportamento estratégico de agente; num
sistema por custo, é prever a saída de um solver.} São problemas estatísticos
diferentes, com fontes de incerteza diferentes, e um modelo treinado no regime
atual não transfere para o seguinte. Isso é risco de obsolescência datada, não
incerteza genérica --- e datar o corte da amostra na transição é obrigatório,
não recomendável.

\section*{Corte de geração: decisão de despacho, não erro de previsão}

É aqui que despacho deixa de ser tema de engenharia e vira o maior problema
econômico do setor elétrico brasileiro. Números do ONS reportados para o
período de janeiro a abril:

\vspace{0.4em}
\tabfont
\begin{tabular}{@{}
    >{\rag\arraybackslash}p{5.0cm}
    >{\centering\arraybackslash}p{2.6cm}
    >{\centering\arraybackslash}p{2.6cm}@{}}
\toprule
\cab{MW médios} & \cab{2025} & \cab{2026} \\
\midrule
Corte total & \num{2436} & \num{2843} \\
\quad como fração do potencial & 15,3\% & 17,2\% \\
\addlinespace[.4ex]
Eólica & \num{1592} & \num{1806} \\
Solar & 844 & \num{1037} \\
\addlinespace[.4ex]
Razão energética & 695 & \textbf{\num{1772}} \\
Confiabilidade elétrica & 591 & 651 \\
Indisponibilidade externa & \textbf{\num{1150}} & 420 \\
\bottomrule
\end{tabular}

\notas{Fonte: ONS, na compilação de~\cite{canalsolar2026}; primeiro
quadrimestre de cada ano. \textbf{Não é medição desta linha} --- é número de
imprensa setorial atribuído ao operador, e não foi conferido contra o dado
aberto do ONS. Distribuição regional em 2026: Nordeste \num{2233} MW, Sudeste
476, Norte 17, Sul 13.}

\vspace{0.6em}
\termo{A inversão da causa dominante é o achado, não o crescimento do total.}
Em 2025 o corte era majoritariamente indisponibilidade externa --- gargalo de
rede, que é problema de obra. Em 2026 é razão energética: a carga não absorve o
que está disponível para despacho. Gargalo se resolve com linha de transmissão;
sobreoferta, não. A concentração de \num{2233} dos \num{2843} MW no Nordeste
confirma que a restrição mudou de natureza sem mudar de lugar.

O próprio ONS enquadra a restrição como medida técnica necessária à segurança
do SIN e não como escolha do operador, e publicou nota técnica sobre critérios
de gestão de excedentes de energia em 23 de setembro de 2025, além de anunciar
em 5 de agosto de 2026 cronograma de atualização do cálculo das
restrições~\cite{onsfaq}. \termo{A métrica está em movimento}, o que é motivo
suficiente para datar qualquer série de corte pela versão do cálculo, e não só
pelo período.

\begin{objecao}
Circula também a cifra de 20,6\% de toda a energia solar e eólica disponível
desperdiçada em 2025, contra os 15,3\% do primeiro quadrimestre do mesmo ano.
Ano inteiro contra quadrimestre, e ``disponível'' contra ``potencial'', são
bases distintas: \textbf{os dois números não são comparáveis, e citá-los lado a
lado produz contradição aparente}. A mesma regra vale para medida própria ---
corte apurado como programado contra previsto no agregado do SIN não é a mesma
grandeza que \emph{constrained-off} apurado por usina contra a geração estimada
dela. Antes de comparar duas frações de corte, casar numerador, denominador e
nível de agregação.
\end{objecao}

Do lado regulatório, a Lei 15.269/2025 é apresentada como base de acerto
administrativo para os cortes ocorridos entre setembro de 2023 e novembro de
2025, com renúncia a ação judicial como contrapartida do ressarcimento; a ANEEL
prorrogou a suspensão do ressarcimento por \emph{constrained-off}; e o MME abriu
a Consulta Pública nº 210, em 31 de dezembro de 2025, sobre as condições de
compensação financeira. Geradores estimam perda da ordem de R\$ 3 bilhões desde
2021~\cite{tagd2026,canalenergia2026}.

Isso importa para modelagem por um motivo específico. \termo{Se o corte passa a
ser compensado, ele deixa de ser custo do gerador e vira custo do consumidor, e
o incentivo de quem opera muda.} Uma série de corte que atravesse a mudança de
regra tem quebra de regime na data da regra, não na física --- e um modelo
ajustado sobre os dois lados da quebra estima uma média de dois regimes que não
descreve nenhum.

\section*{Despacho fora da ordem de mérito}

O outro lado da mesma decisão. Quando o ONS despacha térmica acima da ordem de
mérito por segurança, a diferença vira Encargo de Serviço do Sistema, pago pelo
consumidor, e devido apenas ao gerador térmico que atende à solicitação do
operador~\cite{ess}. Soma-se a isso o deslocamento hidrelétrico causado por
geração térmica inflexível, tratado como custo decorrente de decisão unilateral
do agente térmico e por isso não abatível da indisponibilidade da usina.

Está aqui a assimetria que estrutura o setor: \termo{corte de renovável e
despacho fora do mérito são a mesma decisão vista dos dois lados, e só um dos
dois tem encargo definido há décadas}. A pergunta de pesquisa que essa
assimetria abre --- se o preço sombra de uma restrição de rede deveria remunerar
os dois lados simetricamente --- é a mesma que a regulamentação do
\emph{curtailment} está decidindo por outra via.

\section*{Aprendizado de máquina: três tarefas distintas, frequentemente confundidas}

\campo{Um, acelerar o solver}
É a linha com resultado mais sólido. Prever restrições redundantes, soluções
iniciais viáveis e o subespaço afim onde o ótimo provavelmente cai resolve SCUC
4,3 vezes mais rápido \emph{com} garantia de otimalidade e 10,2 vezes mais
rápido \emph{sem} ela~\cite{xavier2020}. A dificuldade real não é acurácia: é
viabilidade. A saída da rede precisa satisfazer restrição física, e a resposta
da literatura é camada de reparo diferenciável em arquitetura ponta a
ponta~\cite{chen2024}, ou camada que busca viabilidade com
garantia~\cite{nguyen2025}.

\campo{Dois, prever o resultado do despacho}
É a linha mais próxima de dado aberto, e a que dispensa acesso ao solver. O
trabalho de referência é o de Gorka e Roald sobre o CAISO: dois classificadores
por \emph{gradient boosting}, um que identifica corte em tempo real e outro que
prevê presença ou ausência de corte para todos os intervalos de um
dia~\cite{gorka2024}. \termo{Formular corte como classificação binária a partir
do estado do sistema é o desenho canônico}, e quem trabalha nesse alvo cita ou
explica por que difere.

\campo{Três, operar a rede diretamente}
Aprendizado por reforço, no ambiente L2RPN e no simulador Grid2Op, da RTE.
Existe agente vencedor da competição de 2022 e existe esforço de aproximar o
método do fluxo de trabalho real, em planejamento do dia anterior e em controle
em tempo real~\cite{lehna2023}. Não há implantação em operador, e o campo ainda
discute qual é o \emph{benchmark} padronizado~\cite{rl2grid2025} --- o que é
sinal de maturidade baixa, não de falta de esforço.

\vspace{0.6em}
O que as três compartilham: \termo{nenhuma substitui o otimizador, e as que
prometeram isso não chegaram à produção.} O que chegou perto foi o papel de
acelerador com garantia preservada. Confundir as três é o erro de leitura mais
comum da área, porque as três se apresentam sob o mesmo rótulo de ``IA para
despacho'' e respondem a perguntas incompatíveis: uma sobre tempo de solução,
uma sobre previsão, uma sobre controle.

\section*{Emissão marginal: a saída sombra do despacho}

O fator de emissão marginal atribui emissão ao gerador cuja saída muda em
resposta a uma variação de carga --- paradigma consequencial, contra o
atributivo da média do sistema~\cite{emaps2025}. É calculado por modelo de
despacho, combinando ordem de mérito com curva de duração de carga para
identificar qual unidade entra ou sai na margem. Para avaliação de ciclo de vida
consequencial de carga nova recomenda-se o marginal; para inventário, a média.
A fronteira recente é a emissão marginal locacional, formulada como mecanismo de
despacho em duas camadas~\cite{cote2026}.

A consequência prática é uma só, e é de categoria: \termo{fator marginal é
derivado do despacho, logo herda inteiro o erro do modelo de despacho} --- a
variabilidade do dual inclusive. Tratá-lo como grandeza observada, e não como
saída de modelo, é confundir medição com simulação.

\section*{Reserva, inércia e a resposta brasileira}

Com penetração alta de recurso baseado em inversor, o nadir de frequência, a
frequência de acomodação e o tempo de resposta se degradam, e a integração de
renovável reduz a inércia do sistema por diminuir a dependência de gerador
síncrono~\cite{smahi2025}. A regulação europeia já classifica inércia como
serviço ancilar de contratação por procedimento de mercado, e a literatura de
2025 se concentra em dimensionamento de reserva de resposta rápida em sistema de
baixa inércia~\cite{panagi2025}.

A resposta brasileira é o LRCAP 2026, primeiro leilão do país dedicado a
armazenamento: dois certames em dezembro de 2026, um com exigência de conteúdo
nacional e outro aberto a equipamento importado, contratos de 15 anos com
receita fixa corrigida pelo IPCA, suprimento a partir de 1º de agosto de 2028,
potência máxima por 4 horas consecutivas, até 2 ciclos por dia e 366 ciclos
anuais~\cite{epe2026,bess2026}.

O que se contrata é \termo{disponibilidade de potência, não energia produzida}.
Isso é a admissão institucional de que o problema do SIN deixou de ser energia e
passou a ser flexibilidade --- e é a mesma constatação que o giro do corte para
``razão energética'' expressa pelo outro lado. As duas evidências são
independentes: uma vem do desenho de leilão, outra da apuração de restrição.

\section*{O que o campo não resolveu}

\begin{ficha}
\item[Viabilidade] Camada de reparo em \emph{proxy} aprendido funciona
  empiricamente; garantia formal em escala real segue aberta, e é o que separa
  ferramenta de triagem de motor de \emph{clearing}.
\item[Dual do SDDP] Sensibilidade da solução às condições iniciais
  autorregressivas é defeito documentado há quase uma década e afeta
  diretamente qualquer previsão de CMO.
\item[Marginal ou médio] Sem consenso fora da avaliação consequencial estrita,
  e a diferença entre os dois é grande onde a matriz é hidrotérmica.
\item[Oferta e hidro] Se o despacho por oferta melhora ou piora a coordenação
  hidrelétrica não tem experimento: nenhum país com parque hidrelétrico da
  escala brasileira fez essa transição.
\end{ficha}

\section*{Procedência deste levantamento}

\begin{objecao}
Três graus de evidência aparecem acima, e a distinção não é decorativa.
\textbf{Registro bibliográfico conferido}: as onze entradas com DOI foram
verificadas campo a campo contra o Crossref e a API do arXiv em 06/09/2026 ---
autor, container, volume, página, ano. \textbf{Página primária baixada}: a FAQ
de \emph{curtailment} do ONS. \textbf{Resumo de segunda mão}: todo o resto,
inclusive os números de corte de~\cite{canalsolar2026} e as datas regulatórias,
que vieram de resultado de busca e de imprensa setorial, não do ato normativo.
Três lacunas conhecidas: o texto da Lei 15.269/2025 não foi aberto, e duas
fontes a citam para coisas distintas; o artigo de~\cite{gorka2024} não foi lido
por inteiro, porque a ACM devolveu HTTP 403, de modo que método e alvo saem do
registro e não do corpo; e a nota técnica do ONS sobre excedentes é citada, não
lida. \textbf{Nenhum número deste documento foi medido em arquivo}, e por isso
nenhum deles pode sustentar afirmação empírica desta linha sem antes ser
refeito sobre dado próprio.
\end{objecao}

\clearpage
{\sffamily\bfseries\large Referências\par}
\vspace{0.3em}
{\sffamily\fontsize{7.4}{9.4}\selectfont\color{tinta2}
\begin{thebibliography}{99}
\setlength{\itemsep}{0.2em}

\bibitem{xavier2020} Xavier, Á.S., Qiu, F. e Ahmed, S. Learning to solve
  large-scale security-constrained unit commitment problems. \emph{INFORMS
  Journal on Computing}, 2020. \textsc{doi}: 10.1287/ijoc.2020.0976.

\bibitem{chen2022} Chen, W., Park, S., Tanneau, M. e Van Hentenryck, P.
  Learning optimization proxies for large-scale Security-Constrained Economic
  Dispatch. \emph{Electric Power Systems Research}, 213, 108566, 2022.
  \textsc{doi}: 10.1016/j.epsr.2022.108566.

\bibitem{chen2024} Chen, W., Tanneau, M. e Van Hentenryck, P. End-to-end
  feasible optimization proxies for large-scale economic dispatch. \emph{IEEE
  Transactions on Power Systems}, 39(2), 4723--4734, 2024. \textsc{doi}:
  10.1109/TPWRS.2023.3317352.

\bibitem{nguyen2025} Nguyen, H.T. e Donti, P.L. \emph{FSNet: feasibility-seeking
  neural network for constrained optimization with guarantees}.
  arXiv:2506.00362, 31 de maio de 2025.

\bibitem{gorka2024} Gorka, J. e Roald, L. Classification models for forecasting
  and real-time identification of solar curtailment in the California grid.
  \emph{Proceedings of the 15th ACM International Conference on Future and
  Sustainable Energy Systems} (e-Energy '24), 158--169, 2024. \textsc{doi}:
  10.1145/3632775.3661951.

\bibitem{maceira2015} Maceira, M.E.P., Marzano, L.G.B., Penna, D.D.J., Diniz,
  A.L. e Justino, T.C. Application of CVaR risk aversion approach in the
  expansion and operation planning and for setting the spot price in the
  Brazilian hydrothermal interconnected system. \emph{International Journal of
  Electrical Power \& Energy Systems}, 72, 126--135, 2015. \textsc{doi}:
  10.1016/j.ijepes.2015.02.025.

\bibitem{soares2017} Soares, M.P., Street, A. e Valladão, D. On the solution
  variability reduction of stochastic dual dynamic programming applied to
  energy planning. \emph{European Journal of Operational Research}, 258(2),
  743--760, 2017. \textsc{doi}: 10.1016/j.ejor.2016.08.068.

\bibitem{cote2026} Cote, L. e Sun, A. \emph{On locational marginal emissions in
  electricity markets: a two-layered dispatch mechanism and its fundamental
  theorems}. arXiv:2603.19530, 19 de março de 2026.

\bibitem{smahi2025} Smahi, A. e Makhloufi, S. The power grid inertia with high
  renewable energy sources integration: a comprehensive review. \emph{Journal
  of Engineering}, 2025(1), 7975311, 2025. \textsc{doi}: 10.1155/je/7975311.

\bibitem{panagi2025} Panagi, S. e Aristidou, P. Sizing of fast frequency
  response reserves for improving frequency security in low-inertia power
  systems. \emph{Sustainable Energy, Grids and Networks}, 42, 101699, 2025.
  \textsc{doi}: 10.1016/j.segan.2025.101699.

\bibitem{lehna2023} Lehna, M. et al. \emph{Reinforcement learning based power
  grid day-ahead planning and AI-assisted control}. arXiv:2302.07654, 2023.
  Primeiro lugar na competição L2RPN 2022, da RTE. Autoria além do primeiro
  nome não conferida no registro.

\bibitem{rl2grid2025} \emph{RL2Grid: benchmarking reinforcement learning in
  power grid operations}, 2025. Registro não conferido no Crossref; entrada
  vinda de resultado de busca.

\bibitem{cepel2026} CEPEL. \emph{Modelos de otimização energética} ---
  documentação das bibliotecas. Consultado em 06/09/2026.
  \texttt{see.cepel.br/manual/libs}

\bibitem{ccee2026} CCEE. \emph{Modelos computacionais de cálculo do PLD}
  --- NEWAVE, DECOMP e DESSEM. Portal de Aprendizado, consultado em 06/09/2026.

\bibitem{cp218} Ministério de Minas e Energia. Consulta Pública nº 218/2026:
  aperfeiçoamentos do Mercado de Curto Prazo e transição para oferta de
  quantidade de energia. Março de 2026, via resumo secundário.

\bibitem{onsfaq} ONS. \emph{FAQ Curtailment}. Página consultada e baixada em
  06/09/2026. \texttt{ons.org.br/Paginas/faq\_curtailment.aspx}

\bibitem{canalsolar2026} Canal Solar. \emph{Brasil já cortou quase 3 GW médios
  de energia solar e eólica em 2026}, 2026. Números atribuídos ao ONS; não
  conferidos contra o dado aberto do operador.

\bibitem{canalenergia2026} CanalEnergia. \emph{ANEEL prorroga suspensão de
  ressarcimento por constrained-off de eólica e solar}, 2026.

\bibitem{tagd2026} Menezes, M. e Wolf, B. \emph{Quais os avanços regulatórios
  sobre curtailment esperados no Brasil?}, TAGD Advogados, 2026.

\bibitem{ess} \emph{Visão geral do Encargo de Serviço de Sistema (ESS)}.
  Material setorial, consultado em 06/09/2026.

\bibitem{epe2026} EPE. \emph{Leilão de Reserva de Capacidade na Forma de
  Potência 2026}. Página oficial, consultada em 06/09/2026.

\bibitem{bess2026} Brasil BESS. \emph{LRCAP 2026 Armazenamento: guia técnico do
  primeiro leilão de baterias do Brasil}, 2026.

\bibitem{emaps2025} Electricity Maps. \emph{Marginal emissions: what they are,
  and when to use them}. Atualizado em janeiro de 2025.

\end{thebibliography}
}

\end{document}
